\section{Enfoque y materiales}
\subsection{Selección y tratamiento de los casos}
Se realizó una lectura dirigida de documentación pública y de materiales proporcionados por los autores. La selección es intencional: incluye experiencias en las que las prácticas abiertas conectan conocimientos especializados con participación distribuida. Las unidades de análisis son la actividad y sus mecanismos organizativos, no los individuos. Se consideran tres formatos: un hackatón de datos, un editatón bilingüe y trayectorias educativas con datos de física.

El informe de CoAfina 2023 describe una actividad virtual con estudiantes universitarios y tareas de análisis \cite{coafina}. El artículo de Salud para todos aporta un registro publicado de edición colaborativa \cite{salud}. El relato de ATLAS sobre educación en Venezuela documenta el uso de datos abiertos en tesis y acompañamiento académico \cite{atlas}. Estos documentos no constituyen una muestra representativa de América Latina. Su lectura conjunta sirve para identificar mecanismos y formular preguntas que deberán contrastarse en otros contextos.

Para cada caso se distinguen: propósito, forma de participación, producto, apoyo especializado, condiciones de reutilización y evidencia de continuidad. Los números se conservan con su unidad original. Inscripción, edición, equipo y tesis no son categorías intercambiables. No se suman cohortes sin identificar participantes repetidos, ni se interpretan visitas a una página como aprendizaje. El registro de evidencia adjunto identifica la procedencia de cada afirmación y los puntos que requieren revisión de los autores.

\subsection{Alcance de la comparación}
La comparación es analítica, no una clasificación de eficacia. Un hackatón puede priorizar exploración, mientras una tesis exige un compromiso temporal y una evaluación diferentes. La pregunta relevante no es qué formato produce más participantes, sino qué combinación de tareas, apoyo y recursos hace posible una contribución verificable. La cercanía de los autores a las experiencias aporta conocimiento operativo, pero también introduce selección y posibles sesgos favorables; por ello se explicitan los límites de las fuentes.

El presente texto no incorpora entrevistas nuevas, encuestas individuales ni seguimiento longitudinal. Sus propuestas de evaluación son prospectivas. Los resultados publicados de otros trabajos se citan como antecedentes; en particular, el estudio de Salud para todos conserva su identidad y autoría propias. La contribución nueva de esta síntesis es la articulación entre formatos y el diseño de un marco común de documentación y continuidad.
